\chapter{Résultats, interprétation et différenciation}
\label{chap:resultats}

\section{Synthèse des résultats quantitatifs}
Le tableau~\ref{tab:synth} consolide les performances atteintes. Sur le plan du diagnostic,
la plateforme atteint ou dépasse les seuils industriels usuels (exactitude $>$ 95\,\% sur les
cas de roulement, $R^2 = 0{,}90$ sur le pronostic).

\begin{table}[h]
\centering\small
\caption{Synthèse des résultats par tâche et dataset.}
\label{tab:synth}
\begin{tabular}{>{\bfseries}l l l l}
\toprule
Tâche & Dataset & Meilleur modèle & Performance \\
\midrule
Classification & VBL-VA001 & XGBoost & 100{,}0\,\% (4 classes) \\
Classification & CWRU & XGBoost & 97{,}6\,\% (10 classes) \\
Classification & Mechanical Faults & MLP & 89{,}6\,\% (4 classes) \\
Classification & MCC5-THU & CNN 1D & 8 classes d'engrenage \\
Régression RUL & CMAPSS FD001 & XGBoost & MAE 9{,}61 ; $R^2$ 0{,}90 \\
Anomalie (non sup.) & tous & Auto-encodeur + ensemble & seuil P95, vote 2/3 \\
\bottomrule
\end{tabular}
\end{table}

\section{Interprétation physique et industrielle}
Les résultats ne sont pas que des chiffres : ils ont un sens mécanique.
\begin{itemize}[leftmargin=1.5em]
  \item \textbf{Roulements (CWRU).} L'exactitude de 97{,}6\,\% sur 10 classes signifie que la
        plateforme distingue non seulement \emph{où} se situe le défaut (bague interne,
        externe, bille) mais aussi son \emph{intensité}. Les rares confusions concernent des
        sévérités voisines d'un même défaut, ce qui est sans conséquence sur la décision de
        maintenance (le même roulement sera remplacé).
  \item \textbf{Machine tournante (Mechanical Faults).} Le diagnostic à 89{,}6\,\% des
        balourds, désalignements et jeux mécaniques couvre les trois défauts les plus
        fréquents des machines tournantes, validant l'utilité opérationnelle.
  \item \textbf{Pronostic (CMAPSS).} Une erreur moyenne de 9{,}6 cycles permet de planifier
        une intervention avec une marge confortable. Le bon score NASA (12\,249) confirme que
        le modèle évite les sur-estimations dangereuses qui repousseraient à tort
        l'intervention.
  \item \textbf{Anomalie.} La détection non supervisée ajoute un filet de sécurité contre les
        modes de défaillance \emph{inconnus}, absents de tout entraînement supervisé.
\end{itemize}

\section{Capacités validées de bout en bout}
Au-delà des modèles, les fonctionnalités d'industrialisation ont été validées
opérationnellement (tableau~\ref{tab:valid}).

\begin{table}[h]
\centering\small
\caption{Validations fonctionnelles de bout en bout.}
\label{tab:valid}
\begin{tabular}{>{\bfseries}p{4.4cm} p{9.2cm}}
\toprule
Fonctionnalité & Résultat de validation \\
\midrule
Sévérité ISO 10816 & Essai 1\,g/50\,Hz $\rightarrow$ 22{,}3\,mm/s, zone D (conforme au calcul théorique) \\
Baseline par machine & Signal sain $\rightarrow$ anomalie $\approx 5\,\%$ ; dégradé $\rightarrow \approx 96\,\%$ \\
Ingestion + GMAO & Signal critique ingéré $\rightarrow$ ordre de travail P1 créé automatiquement \\
Connecteur OPC-UA & Serveur de démo $\rightarrow$ machine ingérée (états + baseline) \\
Copilot RAG & \og pic à 2$\times$ \fg{} $\rightarrow$ \og désalignement \fg, sources citées \\
KPIs réels & MTBF/MTTR/disponibilité calculés depuis les interventions saisies \\
Lecture agrégée & Centaines de points bruts $\rightarrow$ tranches horaires (downsampling) \\
\bottomrule
\end{tabular}
\end{table}

\section{Apports et différenciation}
La valeur de PrognoSense ne réside pas dans un record ponctuel de précision, mais dans la
\textbf{combinaison cohérente} d'éléments rarement réunis dans un même projet :
\begin{enumerate}[leftmargin=1.6em]
  \item \textbf{Amplitude fonctionnelle.} Une seule plateforme couvre quatre natures de
        machines, cinq familles de défauts (cf. tableau~\ref{tab:fautes}), et trois tâches
        d'apprentissage (classification, régression, détection non supervisée).
  \item \textbf{Rigueur scientifique.} Évaluation sans fuite de données, métriques adaptées
        (score NASA asymétrique), indicateurs de fiabilité \emph{mesurés} et non inventés.
  \item \textbf{Crédibilité normative.} Diagnostic exprimé en zones ISO~10816, architecture
        inspirée d'OSA-CBM (ISO~13374).
  \item \textbf{Actionnabilité.} Boucle fermée de l'alerte vers l'ordre de travail et la
        GMAO --- ce que la plupart des prototypes académiques ne font pas.
  \item \textbf{Confiance et gouvernance.} Explicabilité, calibration, dérive, audit,
        versioning avec rollback, et mesure réelle du taux de fausses alarmes.
  \item \textbf{Ouverture et connectivité.} Ingestion vendor-agnostic, OPC-UA et MQTT,
        absence de verrouillage propriétaire.
  \item \textbf{Accessibilité.} Copilot ancré sur les normes, rendant l'expertise vibratoire
        accessible en langage naturel.
\end{enumerate}
C'est cette \emph{intégration verticale complète} --- du capteur jusqu'à l'ordre de travail,
en passant par le diagnostic normalisé, le pronostic et la gouvernance --- qui démarque
PrognoSense des travaux centrés sur un unique algorithme ou un unique dataset, et qui le
rapproche d'une véritable solution industrielle.
